\section{Introduction}
\label{sec:introduction}

LLM serving systems increasingly execute \emph{workflows}, not isolated
requests.  An agent may invoke several model replicas and tools in parallel,
join their outputs, and issue another round of long-context requests.  The KV
cache accumulated along this execution is expensive to regenerate; after a
worker failure, token replay delays both the interrupted request and downstream
nodes waiting for it.

Existing systems reduce this cost by streaming, checkpointing, or offloading KV
state to peer GPUs, host memory, or storage
~\cite{strati2024dejavu,jayakody2026ghostserve,qin2025mooncake,xie2025hicache}.
They answer how to move and restore KV, but not which live state most reduces
workflow completion delay under scarce protection capacity.  Size, recency,
and reuse cannot distinguish equal-size requests on different parallel
branches, where one recovery pause blocks a join and another is hidden.

A branch's remaining \emph{Slack} is the delay it can absorb before extending
the next mandatory join.  Recovery below that Slack does not change workflow
completion time; delay beyond it is exposed.  The application can derive this
request-level value from visible dependencies, whereas a KV manager cannot.
Slack also changes after every model or tool result, making a static critical
label too coarse.

\system{} applies Slack twice.  At each boundary, a GBDT predicts visible
remaining branch time and CPM computes Slack
~\cite{friedman2001gbdt,kelley1959criticalpath}; the next request carries one
scalar.  During normal operation, the Controller compares the Slack-exposed
harm removed by a complete protected chunk with the harm of evictions.  After
failure, the same Slack orders Peer-HBM adoption, DRAM loading, and
recomputation, aligning protection value and recovery urgency.

Two constraints make this nontrivial.  Slack changes faster than KV should
move, so \textsc{SlackLandlord} updates metadata lazily and each pool globally
sorts only one safe suffix tail per request.  Peer protection must also yield
to inference: \textsc{HBMQuotaAIMD} changes only a target's shadow-HBM quota,
while the Worker allocator remains the capacity authority
~\cite{chiu1989aimd}.

\system{} is designed to handle one fail-stop serving domain at a time and to
recover on a Router-selected, same-layout target using only complete,
version-compatible KV.  Missing or stale state falls back to recomputation;
Slack affects priority, never correctness or capacity.  Our evaluation models
this design with single-worker-shard fail-now counterfactuals rather than
executed worker replacement.

This paper makes four contributions:

\begin{itemize}
  \item It proposes causal workflow Slack as a value signal for KV fault
  tolerance and quantifies the opportunity in a public agent trace
  (\S\ref{sec:background}).

  \item It designs \textsc{SlackLandlord}, which uses one fail-now model for
  independent HBM/DRAM credits and recovery order while maintaining continuous
  prefixes with a two-level victim index (\S\ref{sec:design}).

  \item It separates policy from safety: source-owned paths commit
  independently, \textsc{HBMQuotaAIMD} adapts only the Peer-HBM envelope, and
  Workers enforce capacity and output fencing
  (\S\ref{sec:control-plane}--\S\ref{sec:end-to-end}).

  \item It reports end-to-end WCT from three frozen legacy SGLSim lineages:
  one paired multi-baseline trial, one explicitly labeled DP64 model
  projection, and one heavy-failure logical stress smoke.  The results retain
  unfavorable comparisons and are separated from current-design validation
  (\S\ref{sec:evaluation}).

  \item It separately audits the current design over a frozen 3,110-run
  simulation matrix, reporting pressure--alpha dynamics, zero-alpha stress,
  HBM/DRAM occupancy, reuse-bin admission, copy amplification, and conditioned
  restored-token coverage without presenting those diagnostics as WCT
  (\S\ref{sec:current-mechanisms}).
\end{itemize}
